% Author: Grayson Orr
% Course: IN721: Mobile Application Development

\documentclass{article}
\author{}

\usepackage{graphicx}
\usepackage{wrapfig}
\usepackage{enumerate}
\usepackage{hyperref}
\usepackage[margin = 2.25cm]{geometry}
\usepackage[table]{xcolor}
\usepackage{fancyhdr}
\hypersetup{
  colorlinks = true,
  urlcolor = blue
}
\setlength\parindent{0pt}
\pagestyle{fancy}
\fancyhf{}
\rhead{College of Engineering, Construction \& Living Sciences\\Bachelor of Information Technology}
\lfoot{Practical 02: Kotlin 2\\Version 1, Semester Two, 2021}
\rfoot{\thepage}

\begin{document}

\begin{figure}
    \centering
    \includegraphics[width=50mm]{../../resources/img/logo.png}
\end{figure}

\title{College of Engineering, Construction \& Living Sciences\\Bachelor of Information Technology\\IN721: Mobile Application Development\\Level 7, Credits 15\\\textbf{Practical 02: Kotlin 2}}
\date{}
\maketitle

\section*{Assessment Overview}
In this assessment, you will solve five coding problems using \textbf{Kotlin} in \textbf{IntelliJ IDEA}. This assessment contributes 5\% towards your final mark in \textbf{IN721: Mobile Application Development}.

\section*{Learning Outcomes}
At the successful completion of this course, learners will be able to: 
\begin{enumerate}
	\item Implement \& publish complete, non-trivial, industry-standard mobile applications following sound architectural \& code-quality standards.
	\item Identify relevant use cases for a mobile computing scenario \& incorporate them into an effective user experience design.
	\item Follow industry standard software engineering practice in the design of mobile applications.
\end{enumerate} 

\section*{Assessment Table}
\renewcommand{\arraystretch}{1.5}
\begin{tabular}{|l|l|l|l|l|}
	\hline
	\vtop{\hbox{\strut \textbf{Assessment}}\hbox{\strut \textbf{Activity}}} & \textbf{Weighting} & \vtop{\hbox{\strut \textbf{Learning}}\hbox{\strut \textbf{Outcomes}}} & \vtop{\hbox{\strut \textbf{Assessment}}\hbox{\strut \textbf{Grading Scheme}}} & \vtop{\hbox{\strut \textbf{Completion}}\hbox{\strut \textbf{Requirements}}} \\
	
	\hline
	
	\small Practicals                                                        & \small 10\%        & \small 2, 3                                                           & \small CRA                                                                    & \small Cumulative                                                           \\ \hline
	\small Project                                                          & \small 70\%        & \small 1, 2, 3                                                        & \small CRA                                                                    & \small Cumulative                                                           \\ \hline
	\small Presentation                                                          & \small 20\%        & \small 2, 3                                                        & \small CRA                                                                    & \small Cumulative                                                           \\ \hline
\end{tabular}

\section*{Conditions of Assessment}
You will complete this assessment during your learner managed time, however, there will be availability during the teaching sessions to discuss the requirements \& your progress of this assessment. This assessment will need to be completed by \textbf{Friday, 13 August 2021}. 

\section*{Pass Criteria}
This assessment is criterion-referenced (CRA) with a cumulative pass mark of \textbf{50\%} over all assessments in \textbf{IN721: Mobile Application Development}.

\section*{Authenticity}
All parts of your submitted assessment must be completely your work \& any references must be cited appropriately including, externally-sourced graphic elements. Provide your references in a \textbf{README.md} file. All media must be royalty free (or legally purchased) for educational use. Failure to do this will result in a mark of \textbf{zero} for this assessment.

\section*{Policy on Submissions, Extensions, Resubmissions \& Resits}
The school's process concerning submissions, extensions, resubmissions \& resits complies with \textbf{Otago Polytechnic} policies. Learners can view policies on the \textbf{Otago Polytechnic} website located at \href{https://www.op.ac.nz/about-us/governance-and-management/policies}{https://www.op.ac.nz/about-us/governance-and-management/policies}.

\section*{Submissions}
You must submit all program files via \textbf{GitHub Classroom}. Here is the URL to the repository you will use for your submission – \href{https://classroom.github.com/a/Bvbfy\_\_J}{https://classroom.github.com/a/Bvbfy\_\_J}. Create a new branch called  \textbf{02-kotlin-2} from the \textbf{main} branch by running the command - \textbf{git checkout -b 02-kotlin-2}. This branch will be your development branch for this assessment. Once you have completed this assessment, create a pull request \& assign the \textbf{GitHub} user \textbf{grayson-orr} to a reviewer. \textbf{Do not} merge your own pull request. Late submissions will incur a \textbf{10\% penalty per day}, rolling over at \textbf{5:00 PM}.

\section*{Extensions}
Familiarise yourself with the assessment due date. If you need an extension, contact the course lecturer before the due date. If you require more than a week's extension, a medical certificate or support letter from your manager may be needed.

\section*{Resubmissions}
Learners may be requested to resubmit an assessment following a rework of part/s of the original assessment. Resubmissions are to be completed within a negotiable short time frame \& usually must be completed within the timing of the course to which the assessment relates. Resubmissions will be available to learners who have made a genuine attempt at the first assessment opportunity \& achieved a \textbf{D grade (40-49\%)}. The maximum grade awarded for resubmission will be \textbf{C-}.

\section*{Resits}
Resits \& reassessments are not applicable in \textbf{IN721: Mobile Application Development}.

\section*{Instructions - Learning Outcomes 2, 3}
You have been provided a directory called \textbf{practical-02-kotlin-2} containing five \textbf{Kotlin} files. In these files, write your solutions to the five problems below. 

\subsection*{Problem 1:} 
Write two classes called \textbf{SoftwareDeveloper} \& \textbf{Manager} which inherit from the given \textbf{Employee} class. The \textbf{SoftwareDeveloper} class has one additional class property called \textbf{favProgLang} of type \textbf{String}. The \textbf{Manager} class also has one additional class property called \textbf{employees} of type \textbf{MutableList$<$Employee$>$} \& three functions which add, remove \& display all managed employees. \\

Use the three \textbf{SoftwareDeveloper} objects \& \textbf{Manager} object in the \textbf{main} function to display the expected output.

\begin{verbatim}
  open class Employee(var id: Int, val firstName: String, val lastName: String, val salary: Int) {
      override fun toString() = "${firstName} ${lastName}"
  }

  // Write your SoftwareDeveloper class here

  // Write your Manager class here

  fun main() {
      val sftDevOne = SoftwareDeveloper(1, "Bert", "Watts", 100000, "Cobol")
      val sftDevTwo = SoftwareDeveloper(2, "Sara", "Cain", 75000, "Perl")
      val sftDevThree = SoftwareDeveloper(3, "Samantha", "Baker", 75000, "PHP")
      val manager = Manager(4, "Owen", "James", 150000, mutableListOf(sftDevOne, sftDevTwo))

      // Write your solution here

      // Expected output:
      // Sara Cain
      // Samantha Baker
  }
\end{verbatim}

\subsection*{Problem 2:}
You have been given a class called \textbf{Stack} of type \textbf{String}. Use the \textbf{Stack} object in the \textbf{main} function to display the expected output.

\begin{verbatim}
  class Stack<String>() {
      private val els = mutableListOf<String>()
      fun push(el: String) = els.add(el)
      fun peek(): String = els.last()
      fun pop(): String = els.removeAt(els.size - 1)
      fun isEmpty() = els.isEmpty()
      fun size() = els.size
      override fun toString() = "Stack[${els.joinToString()}]"
  }

  fun main() {
      val stack: Stack<String> = Stack()
      stack.push("Django")
      stack.push("Laravel")
      stack.push("Ruby on Rails")
      stack.push("Spring")

      // Write your solution here

      // Expected output:
      // Stack[Django, Laravel, Ruby on Rails]
      // Ruby on Rails is at the top of the stack
      // There are 3 item(s) in the stack
  }
\end{verbatim}

\subsection*{Problem 3:}
You have been given a class called \textbf{Stack} of type \textbf{String}. Use the \textbf{Stack} object in the \textbf{main} function \& the \textbf{readLine} function to reverse the user's input.

\begin{verbatim}
  class Stack<String>() {
      private val els = mutableListOf<String>()
      fun push(el: String) = els.add(el)
      fun peek(): String = els.last()
      fun pop(): String = els.removeAt(els.size - 1)
      fun isEmpty() = els.isEmpty()
      fun size() = els.size
      override fun toString() = "Stack[${els.joinToString()}]"
  }

  fun main() {
      val stack: Stack<String> = Stack()

      // Write your solution here

      // Expected output:
      // Enter some text: John Doe
      // eoD nhoJ
  }
\end{verbatim}

\subsection*{Problem 4:}
You have been given a class called \textbf{Stack} of type \textbf{Int}. Use the \textbf{Stack} object in the \textbf{main} function \& the \textbf{readLine} function to convert the user's input into binary. 

\begin{verbatim}
  class Stack<Int>() {
      private val els = mutableListOf<Int>()
      fun push(el: Int) = els.add(el)
      fun peek(): Int = els.last()
      fun pop(): Int = els.removeAt(els.size - 1)
      fun isEmpty() = els.isEmpty()
      fun size() = els.size
      override fun toString() = "Stack[${els.joinToString()}]"
  }

  fun main() {
      val stack: Stack<Int> = Stack()

      // Write your solution here

      // Expected output:
      // Enter a number: 50
      // 110010
  }
\end{verbatim}

\subsection*{Problem 5:}
You have been given a class called \textbf{Stack} of type \textbf{Char} \& an incomplete function called \textbf{isBalanced} which accepts a \textbf{String} parameter called \textbf{sequence}. Given a \textbf{sequence} containing only parentheses, curly brackets \& square brackets, determine if \textbf{sequence} is valid.

\begin{verbatim}
  class Stack<Char>() {
      private val els = mutableListOf<Char>()
      fun push(el: Char) = els.add(el)
      fun peek(): Char = els.last()
      fun pop(): Char = els.removeAt(els.size - 1)
      fun isEmpty() = els.isEmpty()
      fun size() = els.size
      override fun toString() = "Stack[${els.joinToString()}]"
  }

  fun isBalanced(sequence: String): Boolean {
      val stack: Stack<Char> = Stack()
      val map = mapOf(
          '(' to ')', ')' to '(',
          '[' to ']', ']' to '[',
          '{' to '}', '}' to '{'
      )
      
      // Write your solution here
  }

  fun main() {
      // Expected output:
      println(isBalanced("{([])}")) // true
      println(isBalanced("{([")) // false
  }
\end{verbatim}

\textbf{sequence} is valid if:
\begin{itemize}
  \item Open bracket must be closed by the same bracket type.
  \item Open bracket must be closed in the correct order.
\end{itemize}

\begin{verbatim}
  // Example 1
  Input: sequence = "()"
  Output: true
  // Example 2
  Input: sequence = "()[]{}"
  Output: true
  // Example 3
  Input: sequence = "{]"
  Output: false
  // Example 4
  Input: sequence = "{[}]"
  Output: false
\end{verbatim}

\end{document}